Intraosseous trancutaneous amputation prostheses (ITAP) are bone-anchored, skin-crossing devices which overcome challenges with socket prostheses and present users with sensory feedback through osseoperception.
Vibrotactile feedback enhances osseoperception and has the possibility to provide touch and grip sensory feedback enabling users to ``feel'' through their prosthesis, however little is understood about the effects of vibration on the long-term health of the device's fixation.
This paper characterises vibrations propagating through ITAP estimated using a computational model, and finds a first mode of vibration at about \SI{930}{Hz}.
At and below clinically relevant frequencies transmission was found to be linear, with a maximum displacement that declines at about \SI{2}{decade\per decade}.
Above \SI{400}{Hz} vibration deviates from linear, displaying a time-variant response and a large change in the relative displacement along the device.